%% ============================================================================
\section{Discussion}
\label{ch:archaeology}
%% ============================================================================

This chapter discusses the archaeological implications and evidential limits of
the engineering results. The thesis concerns how training data
reaches a machine learning model on a shared supercomputer, not how
archaeological sites are identified. That framing is deliberate, but it leaves
a question that a reader from archaeology is entitled to ask directly -- what,
if anything, does this work contribute to archaeological practice, and how far
can its results be trusted?

The discussion answers that question in precise terms. Technical vocabulary is defined where
it is unavoidable, and the boundary between what has been demonstrated here,
what has been demonstrated by others, and what remains untested is stated
explicitly rather than left to inference.

\subsection{Regional transferability}
\label{sec:arch-regions}

\subsubsection{Scope of transfer}
No evidence in this thesis supports direct deployment of a model trained on the
southern Mesopotamian floodplain in a different region. Published work instead
supports adaptation with local examples, including retraining the same underlying
model for the Abu Ghraib district~\cite{Pistola25}.

\subsubsection{Sources of domain shift}
The model does not learn the concept ``tell''. It learns a statistical
association between patterns of colour, tone and texture in an image and the
outlines an archaeologist drew around them. The qualitative analysis in
Section~\ref{sec:repro-qualitative} makes this concrete: predictions follow
tonal and textural contrast, and where a mound is not tonally distinct from the
surrounding land the model misses it. Anything that changes that relationship
between what a site looks like and what the surrounding landscape looks like
degrades performance. Different soil chemistry, vegetation cover, land-use
regime, season of image capture, preservation state, or satellite sensor all
change it.

\subsubsection{Adaptation through transfer learning}
Adapting an existing model to a new region does not mean starting over. The
standard practice, used throughout this thesis, is \emph{transfer learning}: a
model already trained on a very large general image collection is used as a
starting point and then trained further on the specific task, which requires
far fewer examples than training from nothing. The same principle applies one
step further along -- a model trained on Mesopotamian tells is a better
starting point for a new region than a generic one.

The evidence that this works in practice is not from this thesis but from the
research programme the dataset originates from. Pistola et al.~\cite{Pistola25}
retrained the same underlying model on declassified CORONA satellite imagery
from the 1960s for the Abu Ghraib district west of Baghdad -- a different
sensor, a different decade, and a landscape substantially altered by modern
development -- and report a segmentation score above 85\%, together with four
previously unrecorded sites subsequently confirmed by field verification. Two
other groups have applied convolutional networks to CORONA imagery in adjacent
regions: Soroush et al.~\cite{Soroush20} in the Kurdistan Region of Iraq, and
Vadineanu et al.~\cite{Vadineanu24} on mounded settlements in Upper
Mesopotamia.

\subsubsection{Annotation as the binding constraint}
The binding constraint is not computing power and not model design. It is
annotated examples: someone must have already outlined a reasonable number of
known sites in the new region for the model to learn from. In a region with a
well-catalogued survey record this is a matter of data preparation; in a region
without one, the model cannot be adapted at all until that survey work exists,
which is precisely the expensive, slow work the model is meant to assist. This
is a genuine limitation of the whole approach, not of this implementation.

It is worth being clear that no cross-region transfer experiment was run in
this thesis. The claim above rests on the published results cited, not on
evidence collected here.

\subsection{Transferability across archaeological feature classes}
\label{sec:arch-structures}

\subsubsection{Scope across feature classes}
Transfer across feature classes is technically possible because the evaluated
architectures are general-purpose segmentation models rather than algorithms
whose structure encodes mounded settlements~\cite{Fan20,Chen18,Xie21}. Evidence
for performance on a new feature class nevertheless requires new annotations and
evaluation; architectural generality alone is not empirical validation.

The models used here are general-purpose image segmentation networks. They were
not designed for archaeology; they were designed to outline objects in
photographs, and they were adapted to this task purely by being shown examples
of tells. Show them examples of something else and they will learn that
instead. The published record bears this out across markedly different feature
types and sensors: qanat shafts, the surface openings of underground water
channels, in Kurdistan~\cite{Soroush20}; ``princely'' burial mounds in the
Eurasian steppe~\cite{Caspari19}; stone burial mounds (kurgans) in the Altai
Mountains~\cite{Chen21}; and archaeological topography in airborne laser
scanning data, which is not satellite photography at all~\cite{Trier19,Guyot21}.

\subsubsection{Visibility, scale, and difficulty}
Two properties matter more than the choice of model. First, the feature must be
visible in the sensor being used -- a buried structure with no surface
expression will not be found in satellite imagery regardless of the method.
Second, and this thesis provides direct evidence for it, the feature must be
large enough relative to the image tile. Section~\ref{sec:repro-qualitative}
reports that when the annotated site occupies less than roughly 12\% of the
analysed area, the median overlap score across all three architectures is zero:
the models simply fail on small targets. Tells are a comparatively favourable
subject, being large, raised and often tonally distinct. Features that are
small, subtle, or visible only under particular seasonal or lighting conditions
should be expected to perform considerably worse, and the size effect above
should be treated as a warning rather than a footnote.

\subsection{Reliability and interpretation}
\label{sec:arch-reliability}

Reliability depends on distinguishing the aggregate accuracy statistic from its
per-site distribution and on separating repeated model evaluation from the
single-run systems measurements.

\subsubsection{Interpretation of aggregate IoU}
This thesis reproduces the reference study's reported accuracy, reaching
$72.8$--$73.5\%$ against a published $74.17\%$
(Section~\ref{sec:repro-closure}). That figure is an
\emph{intersection-over-union} score: the area where the predicted outline and
the archaeologist's outline agree, divided by the area covered by either of
them. A score of 1.0 is a perfect match and 0.0 is no overlap at all.

It would be natural to read ``73\%'' as meaning that, on a typical site, about
three-quarters of the outline is drawn correctly. That is not what it means,
and the difference is substantial. The score is averaged over all test images,
including images containing no site at all, and an image correctly identified
as containing nothing scores a perfect 1.0. Of the 589 test images, 123 contain
no site. Those images contribute perfect scores to the average without any
outline having been drawn (Section~\ref{sec:repro-metric} quantifies the effect
of this convention).

\subsubsection{Per-site error distribution}
Scoring only the 466 test images that do contain a site gives a much more
uneven result. The median score per image is $0.165$, $0.053$ and $0.000$ for
the three architectures -- that is, on more than half of individual sites, at
least two of the three models produce almost no useful outline. At the same
time, roughly a third of sites score above $0.5$ and the best-performing tenth
score around $0.78$--$0.81$.

The distribution is therefore not centred on a middling value; it is
\emph{bimodal}. The model tends either to find a site and outline it well, or
to miss it almost entirely, with relatively few cases in between. For practical
purposes this is more useful than the average suggests -- a confident, accurate
outline on a third of sites is a real result -- but it is a materially
different claim from ``73\% accurate'', and it should not be reported as the
latter.

\subsubsection{Reliability of systems measurements}
The performance results in Sections~\ref{ch:results} and~\ref{ch:ablation} come
with a limitation stated plainly there: each configuration was run once, not
repeated. On a shared supercomputer, timing measurements vary with whatever
else is running, so single measurements cannot be separated from that
variation. Those results should be read as a broad, consistent survey rather
than as precisely established differences. The accuracy figures are on firmer
ground, having been measured over ten independent evaluation passes with very
close agreement between them.

\subsection{Verification strategy}
\label{sec:arch-verification}

Three checks were applied, and it is important to be clear about what each one
does and does not establish.

\subsubsection{Held-out evaluation}
The models were evaluated on 589 images set aside before training and never
used to adjust the model. This guards against the most basic failure, in which
a model appears accurate merely because it has memorised the examples it was
trained on.

\subsubsection{Repeated evaluation}
Because the evaluation procedure involves a random crop of each image, the
whole test set was evaluated ten times independently. Agreement between those
ten runs was close (a spread of $0.003$--$0.004$), so the reported figure is not
an artefact of one lucky pass.

\subsubsection{Independent reproduction}
The pipeline was rebuilt from scratch and checked against an independently
published result~\cite{Casini23}, landing within $1$--$1.4$ percentage points.
This is the strongest of the three checks, because it tests the whole apparatus
against a number produced by different people using different code.

\subsubsection{Limits of verification}
Every check above compares the model's output against the \emph{annotations} --
the outlines drawn by archaeologists in the source catalogue. It does not
compare them against the ground. A model that agrees perfectly with the
catalogue would score perfectly here even if the catalogue itself were
mistaken, and the catalogue is known to be imperfect in a specific and
consequential way: a substantial number of recorded sites have been levelled
for agriculture or built over and are no longer visible in modern imagery. A
model that correctly reports nothing visible on such a tile is penalised as a
failure.

No site proposed by any model in this thesis has been checked in the field. The
only field verification in this research programme is that reported by Pistola
et al.~\cite{Pistola25}, who confirmed four previously unrecorded sites found
using the CORONA-trained variant of the same model. Field verification is the
only test that can distinguish a genuine discovery from an agreement with an
imperfect catalogue, and it lies outside the scope of this work.

\subsection{Error sources and feature confusion}
\label{sec:arch-confusion}

Yes, in both directions, and this thesis can quantify how often.

\subsubsection{False-positive detections}
Of the 123 test images containing no archaeological site, 108 were correctly
rejected by all three models. Spurious detections appeared on $5.7\%$
(DeepLabV3+), $7.3\%$ (MA-Net) and $8.9\%$ (SegFormer) of them. Since the
models respond to tonal and textural contrast, the features that trigger them
are those that resemble a mound from above without being one: field boundaries
and irrigated parcels, quarries and borrow pits, spoil heaps, patches of bare
or salt-crusted ground, and modern earthworks. These images were themselves
deliberately sampled by the original authors from urbanised areas, intensive
agriculture, flood-prone terrain and rocky hills -- that is, they are hard
cases by construction, and a roughly 6--9\% false-detection rate on them is a
meaningful but not disqualifying figure.

\subsubsection{False-negative detections}
This is the larger error, and it deserves the greater emphasis. Roughly half of
the images that do contain a site yield almost no useful outline
(Section~\ref{sec:arch-reliability}), concentrated overwhelmingly among smaller
sites. A negative result from these models is therefore weak evidence: absence
of a detection is emphatically not evidence of absence of a site.

\subsubsection{Annotation ambiguity}
Some of what is scored as error is disagreement about the annotation rather
than a mistake in the image analysis. Where a catalogued site is no longer
visible, the model and the catalogue disagree and the model is recorded as
wrong. Fiorucci et al.~\cite{Fiorucci22} argue more generally that standard
detection measures are poorly matched to archaeological objects and propose
alternatives; the concern applies directly here.

\subsubsection{Role of expert review}
The reference study addressed exactly this by placing an archaeologist in the
loop to adjudicate the model's proposals, which raised site-level detection
accuracy from roughly 63\% to 80\%~\cite{Casini23}. That is the appropriate
framing for the whole approach, and it is the framing this thesis
inherits: the model is a triage tool that narrows where a specialist looks,
not an authority on what is there. Verschoof-van der Vaart and
Lambers~\cite{Verschoof22} make the same argument from the perspective of
integrating such tools into archaeological practice.

\subsection{Archaeological contribution}
\label{sec:arch-contribution}

Stated conservatively, the contribution is threefold, and only the third is
novel to this work.

First, it independently reproduces a published archaeological detection result.
Reproduction is undervalued: it establishes that the published figure is
attainable by someone other than its authors, and Appendix~\ref{app:defects}
documents the specific defects and environmental sensitivities that had to be
resolved to obtain it, which is information of practical use to anyone
attempting the same.

Second, it characterises the error structure of that result more precisely than
the headline figure allows -- the bimodal distribution, the strong dependence
on site size, and the false-detection rate on hard negatives. This matters for
anyone deciding how much weight to place on such a model's output.

Third, and this is the actual contribution, it reduces the cost of retraining.
The barrier to applying these methods to a new region or a new class of feature
is rarely the model; it is the effort of assembling annotations and the time
and computing resources needed to train on them. This thesis shows that on a
shared supercomputer a substantial fraction of that time is spent not
computing but waiting for image files to be read from storage, and that
addressing this raises GPU utilisation substantially
(Section~\ref{ch:results}). Cheaper retraining does not make the archaeology
more accurate. It makes it more practical to retrain for a new region, a new
sensor, or a new feature type -- which, per
Sections~\ref{sec:arch-regions} and~\ref{sec:arch-structures}, is exactly what
applying this approach anywhere new requires.
